\section*{Chapter \ref{ch:g7:triangles} --- Triangles and Angles}

\begin{solution}{exo:g7:triangles:1}
Complements: $60^\circ$; $45^\circ$; $18^\circ$.
Supplements: $150^\circ$; $135^\circ$; $108^\circ$.
\end{solution}

\begin{solution}{exo:g7:triangles:2}
The \omterm{def:g7:triangles:pairs}{vertically opposite} angle also measures $35^\circ$
(\cref{thm:g7:triangles:equalities}); the two remaining angles are
supplementary to it: $180 - 35 = 145^\circ$ each.
\end{solution}

\begin{solution}{exo:g7:triangles:3}
(a) $180 - 100 = 80^\circ$. \quad
(b) $180 - 118 = 62^\circ$. \quad
(c) $180 - 150 = 30^\circ$.
\end{solution}

\begin{solution}{exo:g7:triangles:4}
Apex $40^\circ$: the base angles share $180 - 40 = 140^\circ$
equally: $70^\circ$ each.

Base angle $70^\circ$: the two base angles are equal, so the apex is
$180 - 2 \times 70 = 40^\circ$.
\end{solution}

\begin{solution}{exo:g7:triangles:5}
$(6, 8, 12)$: $12 < 6 + 8 = 14$: triangle exists.

$(5, 5, 11)$: $11 > 5 + 5 = 10$: impossible.

$(4, 9, 13)$: $13 = 4 + 9$: flat ``triangle'' --- the points are
aligned, no genuine triangle.

$(7, 7, 7)$: $7 < 14$: \omterm{def:g6:shapes:triangles}{equilateral triangle}.
\end{solution}

\begin{solution}{exo:g7:triangles:6}
$x + 2x + 3x = 6x = 180^\circ$, so $x = 30^\circ$: the angles are
$30^\circ$, $60^\circ$ and $90^\circ$ --- a \omterm{def:g6:shapes:triangles}{right triangle}.
\end{solution}

\begin{solution}{exo:g7:triangles:7}
Prediction: $\widehat{BAC} = 180 - (50 + 60) = 70^\circ$; the
protractor on the finished figure confirms it.
\end{solution}

\begin{solution}{exo:g7:triangles:8}
At the first crossing, the four angles are $54^\circ$, $126^\circ$,
$54^\circ$, $126^\circ$ (\omterm{def:g7:triangles:pairs}{vertically opposite} pairs and supplements).
The \omterm{def:g6:lines:perp}{parallel} line reproduces the same four measures at the second
crossing (corresponding angles): eight angles in all, four of
$54^\circ$ and four of $126^\circ$.
\end{solution}

\begin{solution}{exo:g7:triangles:9}
Let $c$ be the third side. The triangle inequality demands
$c < 8 + 3 = 11$ and also $8 < c + 3$, i.e.\ $c > 5$. So
$5 < c < 11$: the length $7$ cm works; the length $12$ cm (or $4$ cm)
does not.
\end{solution}

\begin{solution}{exo:g7:triangles:10}
$\widehat A + \widehat B = 180 - 90 = 90^\circ$, and
$\widehat A = 2\widehat B$, so $3\widehat B = 90^\circ$:
$\widehat B = 30^\circ$ and $\widehat A = 60^\circ$.
\end{solution}

\begin{solution}{exo:g7:triangles:11}
The diagonal $[AC]$ splits $ABCD$ into triangles $ABC$ and $ACD$.
The four angles of the \omterm{def:g4:geometry:polygon}{quadrilateral} are exactly the six angles of
the two triangles, regrouped (the angles at $A$ and at $C$ are each
split in two). Total: $180 + 180 = 360^\circ$. A \omterm{def:g4:geometry:polygon}{pentagon} splits from
one \omterm{def:g5:solids:def}{vertex} into three triangles: $3 \times 180 = 540^\circ$.
\end{solution}
